\chapter{Úvod}
Hlavním cílem této bakalářské práce bylo vytvoření komponenty pro výukový server v oblasti teoretické informatiky. V praktické části bylo potřeba navrhnout a implementovat dynamickou webovou aplikaci, která by sloužila jako názorný a interaktivní nástroj pro pochopení vybraných algoritmů a principů amortizované složitosti. Výsledná aplikace má studentům umožnit sledovat průběh jednotlivých operací, změny stavů datových struktur a další důležité hodnoty potřebné pro lepší pochopení chování algoritmů i jejich amortizované složitosti.

Problematika amortizované složitosti patří mezi velmi důležitá témata při studiu algoritmů a datových struktur. Její význam spočívá v tom, že umožňuje lépe popsat chování algoritmů v situacích, kdy jednotlivé operace nemají vždy stejnou cenu a jejich skutečná náročnost se projeví až v rámci delší posloupnosti kroků. U některých algoritmů a datových struktur nestačí hodnotit jednotlivé operace samostatně, protože jejich skutečná efektivita se ukáže teprve při sledování delší posloupnosti operací.

I když je amortizovaná složitost významnou součástí TI, bývá pro mnoho studentů obtížně pochopitelná. Hlavním důvodem je její abstraktnější charakter, protože při snaze o její pochopení nestačí sledovat pouze výsledek jedné operace, ale je také nutné porozumět souvislostem mezi více operacemi, změnám stavů datových struktur a principům, podle nichž se náklady mezi jednotlivými kroky rozkládají. Z tohoto důvodu byla tato bakalářská práce vytvořena jako podpora teoretického výkladu prostřednictvím interaktivní vizualizace a možnosti sledovat chování algoritmů na různých příkladech.

Výsledkem této práce je dynamická webová aplikace, která může sloužit budoucím studentům jako výuková pomůcka pro názornou vizualizaci vybraných algoritmů. Důraz je kladen především na možnost aktivní práce uživatele se vstupy a na sledování chování algoritmů v jednotlivých krocích. Student tak může lépe porozumět tomu, jak se mění stav datové struktury, jak se vyvíjí průběh operací a proč nelze u některých algoritmů posuzovat jejich efektivitu pouze na základě jedné samostatné operace. Výsledná komponenta má za cíl propojit teoretický výklad s praktickou ukázkou a přispět k názornějšímu pochopení celé problematiky.

Tato práce je rozdělena do několika hlavních kapitol. Po úvodní kapitole následuje kapitola věnovaná základní problematice amortizované složitosti, ve které jsou vysvětleny klíčové pojmy a principy související s touto oblastí. Další kapitola se zaměřuje na návrh komponenty výukového serveru a na popis jejích hlavních vlastností a požadavků. Následující kapitola se věnuje samotné implementaci aplikace a použitým technologiím. Závěrečná kapitola shrnuje dosažené výsledky a zhodnocuje přínos vytvořeného řešení.